\section{Matchings}

\begin{definition}
    An edge set $M \subseteq E$ is called a \textbf{matching} in $G = (V, E)$ if no vertex is incident to more than one edge from $M$.
\end{definition}

\begin{definition}
    A vertex $v$ is \textbf{covered} by $M$ if it is incident to an edge in $M$.
\end{definition}

\begin{definition}
    A matching $M$ is
    \begin{itemize}
        \item \textbf{perfect} if every vertex is covered ($|M| = |V|/2$).
        \item \textbf{maximal} if it cannot be extended by adding edges.
        \item \textbf{maximum} if there is no matching with more edges.
    \end{itemize}
\end{definition}

\begin{theorem}
    The geedy algorithm leads to a maximal matching $M_{greedy}$ with $\mid M_{greedy} \mid \geq \frac{1}{2} \mid M_{max} \mid$.
\end{theorem}

\begin{proof}
    Observe that every edge in $M_{max}$ covers at least one edge of $M_{greedy}$ else we would be able to add it to $M_{greedy}$. Further more every edge in $M_{greedy}$ covers at most two edges of $M_{max}$.
\end{proof}

\begin{algorithm}
M = empty set
while E is not empty:
    choose an arbitrary edge e from E
    add e to M
    remove all edges incident to the endpoints of e from E
\end{algorithm}

\begin{definition}
    An \textbf{$M$-augmenting path} is a path that alternates between edges in $M$ and not in $M$, starting and ending in uncovered vertices.
\end{definition}

\begin{remark}
    If we fing an $M_{greedy}$-augmenting path we can switch the edges along that path to increase the size of the matching by 1.
\end{remark}

\begin{theorem}
    \textbf{Hall's Marriage Theorem}: A bipartite graph $G = (A \cup B, E)$ contains a matching of cardinality $|M|=|A|$ iff:
    $$\forall X \subseteq A : |X| \le |N(X)|$$
\end{theorem}

\begin{proof}
    $\\\implies$
    This is trivial as if such a matchin exists then of course for every subset $X \subseteq A$ we have $|X| \le |N(X)|$ as we can  match every vertex in $X$ with a distinct vertex in $N(X)$.
    $\\\impliedby$
    We show this using induction over the size of $A$. The BC is $\mid A \mid = 1$ is trivially true. 
    If $\forall |X| \le |N(X)|$ holds then we chose a random edge $(u, v)$  and delete all incident edges and show that the theorem still holds. 
    If $\exists |X_0| = |N(X_0)|$ then we show that the graphs $G_1 = (A\setminus X_0 \cup B\setminus N(X_0))$ and $G_2 = (X_0 \cup N(X_0))$ the theorem still holds.
\end{proof}

\begin{theorem}
    \textbf{Frobenius Corollary}: Every $k$-regular bipartite graph contains a perfect matching.
\end{theorem}

\begin{remark}
    It even holds that every $k$-regular bipartite graph is a union of perfect matchings.
\end{remark}

\begin{theorem}
    In a $k$- and $2^k$-regular bipartite graph we can find a perfect matching in $O(|E|)$ time.
\end{theorem}

\begin{theorem}
    In bipartite grapahs we can find a perfect mathcin in time $O(|V| |E|)$.
\end{theorem}

\begin{remark}
    Important \textbf{matching-algorithms}:
    \begin{itemize}
        \item $O({|V|}^{\frac{1}{2}}|E|)$ Hopcroft-Karp (bipartite unweighted)
        \item $O({|E|}^{1+O(1)})$ Hpcroft-Karp (bipartite polynomial-weighted)
        \item $O({|V|}^{\frac{1}{2}}|E|)$ Micali-Vazirani (general unweighted)
        \item $O({|V|}^{\frac{1}{2}}|E|)$ Gabow-Tarjan (general polynomial-weighted)
        \item $O({|V|}^{2.373})$ Mucha-Sankowski (general unweighted)
    \end{itemize}
\end{remark}

\begin{theorem}
    \textbf{Berge's Theorem (1957)}: A matching $M$ is maximum iff there is no $M$-augmenting path.
\end{theorem}

\begin{proof}
    Else we change the edges along the augmenting path to increase the size of the matching by 1 which contradicts that it was a maximum matching.
\end{proof}

\begin{algorithm}
#BFS to find augmenting paths in bipartite graphs

L(0) = {uncovered vertices in A} (mark all as visited)

for i: 1 - n:
    if i odd:
        L(i) = {unvisited neighbours of L(i-1) via edges in E \ M}
    else:
        L(i) = {unvisited partners of L(i-1) via edges in M}
    Mark L(i) as visited
    
    if vertex v in L(i) not covered 
        return path to v
\end{algorithm}

\begin{algorithm}
#Hopcroft-Karp Algorithm

while there exists an augmenting path P:
    find a maximal set of vertex-disjoint shortest augmenting paths
    augment M along all these paths
\end{algorithm}

